\section{Institutional decomposition}
\label{sec:decomposition}

The composite index may conceal relationships concentrated in one institution. \Cref{tab:decomp}
reports within-standardised M3 coefficients and primary Granger-style joint tests. Positive values are
oriented toward more corruption for every outcome, and FDR correction covers the four subtype
association tests and eight subtype-by-direction Granger tests.

\begin{table}[H]
\centering
\caption{Primary corruption-dimension results.}
\label{tab:decomp}
\begin{threeparttable}
\small
\begin{tabular}{lrrrrr}
\toprule
& \multicolumn{2}{c}{Contemporaneous M3} & \multicolumn{2}{c}{Granger joint $p$} \\
\cmidrule(lr){2-3}\cmidrule(lr){4-5}
Outcome & Std. $\beta$ & FDR $q$ & $P\rightarrow C$ & $C\rightarrow P$ & $N_{M3}$ \\
\midrule
Composite benchmark & -0.097 & -- & 0.686 & 0.241 & 3,940 \\
Executive & 0.021 & 0.685 & 0.496 & 0.870 & 3,940 \\
Public sector & -0.097 & 0.113 & 0.734 & 0.227 & 3,940 \\
Legislative & -0.076 & 0.213 & \textbf{0.003} & 0.064 & 3,745 \\
Judicial & \textbf{-0.173} & \textbf{0.009} & 0.316 & 0.251 & 3,940 \\
\bottomrule
\end{tabular}
\begin{tablenotes}[flushleft]\footnotesize
\item Granger cells show raw joint $p$ values for compactness. The legislative forward-direction test
has FDR $q=0.028$; all other subtype direction tests have $q>0.25$. The composite is a benchmark and is
excluded from subtype FDR families.
\end{tablenotes}
\end{threeparttable}
\end{table}

\subsection{Contemporaneous dimension results}

The negative composite association is not uniform. Executive corruption is near zero
($\widetilde{\beta}=0.021$, $q=0.685$). Public-sector and legislative corruption are negative but
imprecise after FDR correction. Judicial corruption has the clearest negative association
($\widetilde{\beta}=-0.173$, 95\% CI $[-0.284,-0.063]$, raw $p=0.002$, $q=0.009$).

\begin{figure}[H]
  \centering
  \includegraphics[width=0.82\textwidth]{decomposition_twfe.png}
  \caption{Controlled two-way fixed-effects association by corruption dimension. Effects are scaled by
  within-country standard deviations so that magnitudes are comparable.}
  \label{fig:decomposition}
\end{figure}

The common complete-case sample preserves the ranking. The judicial estimate remains negative
($\widetilde{\beta}=-0.144$, $q=0.042$), so the result is not explained by its larger available sample.
Its point estimate is unchanged under country, country-and-year, and Driscoll--Kraay covariance
estimators. A contemporaneous first-difference model retains the negative sign but does not survive
four-outcome correction (raw $p=0.033$, $q=0.134$). The coarse five-level judicial measure is also
negative ($\widetilde{\beta}=-0.137$, $p=0.026$). The sign is therefore robust to measurement and
inference choices, while confirmatory significance is weaker under differencing.

The judicial interaction mirrors the composite: the populism slope is negative at low and median
polyarchy and approximately zero at high polyarchy. This secondary result is consistent with
institutional heterogeneity, but it may also reflect low within-country resolution or expert coding in
less democratic regimes.

\subsection{Dimension-specific temporal ordering}

No subtype Granger test is significant after FDR except populism to legislative corruption (raw joint
$p=0.0035$, $q=0.0279$). The lag-1 coefficient is $-0.0328$ ($p=0.236$); the lag-2 coefficient is
$+0.0733$ (95\% CI $[0.0234,0.1231]$, $p=0.004$). In within-standardised terms, the two-year effect is
only $+0.0234$. The reverse legislative direction does not survive multiplicity correction
($p=0.064$, $q=0.257$).

\begin{figure}[H]
  \centering
  \includegraphics[width=0.74\textwidth]{decomposition_granger_fdr.png}
  \caption{FDR-adjusted joint Granger-style tests across corruption dimensions and directions. Only the
  populism-to-legislative-corruption family falls below 0.05.}
  \label{fig:decompgranger}
\end{figure}

The legislative signal passes the common-sample and placebo-lead checks: future populism does not
jointly predict current legislative corruption after FDR ($q=0.239$). However, it is not confirmed by
first differences (raw $p=0.374$), election-year observations ($p=0.701$), or government-change years
($p=0.486$). Simpler lead-lag models retain a signal in the full and forward-fill-restricted samples.

The evidence is consequently specific but fragile. A positive two-year coefficient is compatible with
slower legislative capture, but dependence on level-based, forward-filled information prevents a
causal interpretation. The decomposition refines the composite null result; it does not overturn it.
